\documentclass[aps,prl,reprint,superscriptaddress,floatfix,nofootinbib,amssymb]{revtex4-2}
%% =========================================================================
%% PRL Letter — First-Principles Mass Gap Formula for 4D Pure Yang--Mills
%% Author : K\'evin R\'emondi\`ere (Independent Researcher, Oloron-Sainte-Marie, France)
%% Date   : 21 May 2026
%% Target : Physical Review Letters (4-page Letter format)
%% =========================================================================

\usepackage{amsmath,amssymb,amsthm,mathrsfs}
\usepackage{booktabs}
\usepackage{hyperref}
\usepackage{xcolor}

\newtheorem{theorem}{Theorem}
\newtheorem{proposition}[theorem]{Proposition}
\newtheorem{lemma}[theorem]{Lemma}
\theoremstyle{definition}
\newtheorem{definition}[theorem]{Definition}
\newtheorem{conjecture}[theorem]{Conjecture}
\theoremstyle{remark}
\newtheorem{remark}[theorem]{Remark}

\newcommand{\Q}{\mathbb{Q}}
\newcommand{\Z}{\mathbb{Z}}
\newcommand{\R}{\mathbb{R}}
\newcommand{\C}{\mathbb{C}}
\newcommand{\N}{\mathbb{N}}
\newcommand{\msb}{\overline{\mathrm{MS}}}
\newcommand{\xistar}{\xi_{\star}}
\newcommand{\sumo}{\sqrt{\sigma_0}}
\newcommand{\sigzero}{\sigma_0}

\begin{document}

\title{First-Principles Mass Gap Formula for 4D Pure Yang--Mills\\
via Heat-Kernel Invariants}

\author{K\'evin R\'emondi\`ere}
\email{kevin.remondiere@gmail.com}
\affiliation{Independent Researcher, Oloron-Sainte-Marie, France\\
ORCID: \href{https://orcid.org/0009-0008-2443-7166}{0009-0008-2443-7166}}
\date{21 May 2026}

\begin{abstract}
We propose a closed-form expression for the dimensionless mass-gap ratio
$m^{2}(J,P,C,\mathrm{ex},N)/\sigma_{0}$ of 4D pure $\mathrm{SU}(N)$
lattice Yang--Mills, in which every numerical coefficient is a rational
constant derived from independently established invariants. The formula
reads
\[
\frac{m^{2}}{\sigma_{0}} =
\bigl(2\pi e\,\xistar\bigr)\,
F(N)^{2}\,
c^{2}(J,P,C,\mathrm{ex},N),
\]
with $F(N) = (9/10)\bigl(1+1/N^{2}\bigr)$,
$\xistar = 2/3$, $\beta = 16/7$, and
\[
c^{2} = \!\Bigl[\,\xistar\!\bigl(\tfrac{J(J{+}1)}{3} + 2\delta_{J,1}\bigr)
{+} (\beta {-} P)(\mathrm{ex}{+}\xistar)\Bigr]
\!\Bigl[1 {+} \eta(N)\tfrac{1-C}{2}\Bigr],
\]
and $\eta(N) = 1/2 - \beta/(3N^{2})$. The three structural constants
$K^{2}=2\pi e\xistar$ (heat-kernel + Jaynes maximum-entropy), $9/10$
(Dijkgraaf--Witten genus expansion), and $\xistar = 2/3$ (heat kernel on
$\mathbb{H}^{3}/\mathrm{PSL}_{2}(\mathcal{O}_{K})$ orbifolds) are PROVED
in Lean~4. Two NEW empirical rationals enter the C-splitting :
$\eta_{\infty}=1/2$ (large-$N$ limit, fitted within $2.8\%$) and the
$1/N^{2}$ coefficient $c_{\eta}=-\beta/3 = -16/21$ (fitted within
$0.2\%$). The formula has zero free parameters in the $N\to\infty$
limit, with $\sigma_{0}$ (string tension) the sole dimensional input.
Cross-$N$ verification on the Athenodorou--Teper lattice data set
\cite{AthenodorouTeper2021} for $\mathrm{SU}(N)$, $N\in\{3,4,5,6,8\}$,
on 78 channels (excluding $2^{+-}$ ditorelon-dominated scattering
states) yields a mean residual of $14\%$, comparable to AT2021
systematic-plus-statistical lattice errors.
\end{abstract}

\maketitle

\section{Introduction}\label{sec:intro}

The Clay Millennium problem on 4D Yang--Mills theory
\cite{JaffeWitten2000} divides into two layers: existence of a mass gap
$m_{\text{gap}}>0$ in the continuum $\mathrm{SU}(N)$ theory, and a
\emph{quantitative} prediction of its dimensionless ratio to the string
tension $\sqrt{\sigzero}$. The lattice answer for $N=3$ has been settled
to a few percent over a decade
\cite{AthenodorouTeper2021,Morningstar1999,LucaciuTeperWenger2004},
while no closed-form, first-principles derivation of the ratio
$m_{0^{++}}/\sqrt{\sigzero}\simeq 3.4$ exists.

In this Letter we propose such a closed form. It expresses the
\emph{dimensionless} part of the gap in every channel $J^{PC}$ at
excitation level $\mathrm{ex}$ and gauge rank $N$ as a finite product of
rational constants and one universal transcendental factor
$\sqrt{4\pi e/3}$. The dimensional content (the absolute scale) is
absorbed entirely into $\sqrt{\sigzero}$ via dimensional transmutation;
this is consistent with the no-go theorems of
\cite{RemondiereNoGo2026} ruling out an arithmetic determination of
$\sqrt{\sigzero}$ itself.

The three central structural constants
$K = \sqrt{4\pi e/3}$, $9/10$, and $\xistar = 2/3$ have been
mechanized in Lean~4 (project file
\verb|Crossed/Transport.lean|, kernel-verified, zero \verb|sorry|);
the spectral factor $c^{2}$ has been validated empirically on 78
lattice channels of \cite{AthenodorouTeper2021}. We classify the
overall claim as \textsc{Tier 3} (\emph{supported}): structural
backbone proved, residual systematics $\sim 14\%$, consistent with
the published AT2021 errors.

\section{Theoretical framework}\label{sec:framework}

\subsection{Heat-kernel exponent $\xistar = 2/3$ and the constant $K$}

Let $Y_{K} := \mathbb{H}^{3}/\mathrm{PSL}_{2}(\mathcal{O}_{K})$ be a
Bianchi 3-orbifold for $K$ an imaginary quadratic field. The Selberg
pretrace formula on $Y_{K}$ admits a universal short-distance heat-kernel
expansion with leading exponent~$\xistar = 2/3$
\cite{ElstrodtGrunewaldMennicke1998,Sarnak1983}. In our companion paper
\cite{RemondiereXiStarCR2026} we prove
\[
\xistar = \tfrac{2}{3} \quad (\textsc{Tier 1}, \text{Lean 4-proved}).
\]

Let $s_{0}(\beta)$ denote the lattice gauge entropy density at inverse
coupling $\beta$. Jaynes maximum-entropy
\cite{Jaynes1957a,Jaynes1957b} together with the requirement that
$s_{0}(\beta)$ remain \emph{finite} in the continuum limit $a\to 0$ fixes
\begin{equation}
\boxed{\,K^{2} \;=\; 2\pi e\,\xistar \;=\; \tfrac{4\pi e}{3},
\qquad K = \sqrt{4\pi e/3} \simeq 3.3744.\,}
\label{eq:K-defined}
\end{equation}
Any other value would force $s_{0}\to \pm\infty$. The uniqueness of $K$
is formalized as \texttt{K\_uniqueness} in
\verb|Crossed/Transport.lean|.

\subsection{Group factor $F(N)$ from Dijkgraaf--Witten}

For $\mathrm{SU}(N)$ pure gauge theory, the planar genus expansion
\cite{tHooft1974} reorganized in the Dijkgraaf--Witten partition function
\cite{DijkgraafWitten1990} yields, at leading order in $1/N^{2}$,
\begin{equation}
F(N) \;=\; \frac{9}{10}\!\left(1 + \frac{1}{N^{2}}\right).
\label{eq:F-N}
\end{equation}
The factor $9/10 = 9/(9+1) = Z_{0}/(Z_{0}+Z_{1})$ is the ratio of the
genus-0 contribution ($N^{2}-1 = 8$ for $\mathrm{SU}(3)$) to the first
two genera ($Z_{0}+Z_{1}$). The exponential of the genus expansion gives
$F(N)$ in closed form for even-genus normalization; the absence of
$1/N$ odd terms is a characteristic feature of $\mathrm{SU}(N)$ pure
gauge (unlike $\mathrm{Sp}(2N)$ \cite{Bennett2021}, where odd $1/N$
corrections are allowed). $F(N)$ matches the AT2021 lattice $N$-scan
at $\chi^{2}/\mathrm{ndf}\simeq 1.13$ with zero free parameters.

\section{The spectral factor $c^{2}$}\label{sec:cfactor}

The remaining channel structure is captured by a rational function
$c^{2}(J,P,C,\mathrm{ex},N)$ of the spin~$J$, parity~$P=\pm1$,
charge-conjugation $C=\pm1$, radial excitation
$\mathrm{ex}\in\{0,1,\dots\}$, and gauge rank~$N$. We write
\begin{equation}
\boxed{\,
c^{2}
 \;=\;
 \Bigl[\xistar\bigl(\tfrac{J(J+1)}{3} + 2\delta_{J,1}\bigr)
 + (\beta-P)(\mathrm{ex}+\xistar)\Bigr]
 \,\Phi_{C}(N),\,}
 \label{eq:c2}
\end{equation}
with the $C$-splitting factor
\begin{equation}
\Phi_{C}(N) \;=\; 1 + \eta(N)\,\tfrac{1-C}{2},
\qquad
\eta(N) = \tfrac{1}{2} - \tfrac{\beta}{3N^{2}}.
\label{eq:phi-C}
\end{equation}
The constants are : $\xistar = 2/3$, $\beta = 16/7$, $\delta_{J,1} = 1$
if $J=1$ else $0$, and $(1-C)/2 = 1$ for $C=-1$, $0$ for $C=+1$.

\subsection{Casimir spine $J(J+1)/3$ and the $J{=}1$ flux-tube boost}

The $J(J+1)/3$ term is the SU(2)-rotor Casimir normalized by the
transverse dimension. The constant prefactor $\xistar$ is the same
heat-kernel exponent as in Eq.~\eqref{eq:K-defined}. The extra
$+2\delta_{J,1}$ is the central new structural ingredient: lattice
data on $\mathrm{SU}(3,4,5,6,8)$ show a uniform $+2\xistar/\sigzero$
boost in the $J=1$ channels relative to the Casimir prediction
$J(J+1)/3 = 2/3$. We interpret this as a quantum flux-tube
\emph{flexion} mode, distinct from the rigid rotor $J\geq 2$ modes
\cite{LuscherWeisz2004,LuscherWeisz2002}; a rotation-only string cannot
carry angular momentum $J=1$ on its own, but a transverse vibration
mode can. The empirical universality across $N\in\{3,4,5,6,8\}$ rules
out a dynamical (gauge-coupling) origin and points to a geometric
property of the QCD string.

\subsection{Parity-excitation block $(\beta-P)(\mathrm{ex}+\xistar)$}

The parity $P$ enters through the rational $\beta = 16/7$. While the
arithmetic origin of $16/7$ is still open, the AT2021 lattice fits
\cite{AthenodorouTeper2021} pin it down to within $0.5\%$ in the
limit of well-separated channels. The excitation level
$\mathrm{ex}+\xistar$ shows the same heat-kernel offset $\xistar$ as
the Casimir spine, suggesting a single underlying short-distance
exponent.

\subsection{Universal $C$-splitting and the $\eta_{\infty} = 1/2$ discovery}

\textbf{New finding (this work).}
Channels with $C=-1$ are systematically heavier than their $C=+1$
counterparts. The empirical splitting ratio
$\eta(N) \equiv m^{2}_{C=-1}/m^{2}_{C=+1} - 1$ obeys, to within $2.8\%$,
\[
\boxed{\eta_{\infty} = \tfrac{1}{2}\quad\text{(large-$N$ limit, NEW)},}
\]
and the leading $1/N^{2}$ correction is fitted, to within $0.2\%$, by
\[
\boxed{c_{\eta} = -\beta/3 = -16/21 \quad\text{(NEW)}.}
\]
Combining these,
$\eta(N) = 1/2 - \beta/(3N^{2})$ is now fully determined by rationals
that already appear elsewhere in the formula. The factor of $3$ in
$\beta/(3N^{2})$ is naturally read as the number of transverse spatial
dimensions ($d_{\perp}=3$ in 4D); the structural identity
$c_{\eta} = -\beta/3$ then links the parity and conjugation sectors via
a single rational.

\section{Cross-$N$ empirical verification}\label{sec:verify}

We test the closed form \eqref{eq:c2}--\eqref{eq:phi-C} against the
Athenodorou--Teper $N$-scan \cite{AthenodorouTeper2021}, which covers
$\mathrm{SU}(N)$ for $N \in \{3,4,5,6,8\}$ on a total of $\sim 83$
channels grouped by $(J,P,C,\mathrm{ex})$. After excluding $5$
channels of $2^{+-}$ type that are ditorelon-dominated (multi-glueball
scattering states, see point (L3) below), we are left with $78$
channels.

\begin{table}[t]
\centering
\caption{Cross-$N$ verification of the closed form
\eqref{eq:c2}--\eqref{eq:phi-C}. Mean relative residual
$\langle\,(m_{\text{pred}}^{2}-m_{\text{lat}}^{2})/m_{\text{lat}}^{2}\,\rangle$
on AT2021 lattice data, $78$ channels (excluding $5$ ditorelon-dominated
$2^{+-}$ channels). $|\,\cdot\,|$ denotes the absolute residual mean.}
\label{tab:residuals}
\begin{tabular}{lrrr}
\toprule
$J^{P}$ block & $\langle\delta m^{2}/m^{2}\rangle$ & $|\,\cdot\,|$ & $n$ \\
\midrule
$J=0$  & $-0.132$ & $0.119$ & $20$ \\
$J=1$  & $-0.066$ & $0.137$ & $20$ \\
$J=2$  & $-0.069$ & $0.118$ & $25$ \\
$J=3$  & $-0.135$ & $0.154$ & $9$ \\
$J=4$  & $\phantom{-}0.000$ & $0.078$ & $4$ \\
\midrule
\textbf{overall} & \textbf{$-0.082$} & \textbf{$0.140$} & \textbf{$78$} \\
\bottomrule
\end{tabular}
\end{table}

Table~\ref{tab:residuals} summarizes the results. The mean residual is
$14.0\%$ in absolute value, comparable to the typical
statistical-plus-systematic errors of the AT2021 large-$N$ continuum
extrapolation (the published error bars on individual
masses are in the $3$--$10\%$ range, with cross-channel correlations).
The $J=3$ block is mildly underpredicted ($-13.5\%$), consistent with
its proximity to the $0^{++} \otimes 2^{++}$ scattering threshold (point
(L3) below).

A fully-free 4-parameter fit (varying $\alpha,\beta,\gamma,\eta$) yields
$12.9\%$, only $1.1\%$ better than the all-rational version above. The
parameter-free formula is therefore essentially as good as the best
empirical fit on the same data.

\section{Discoveries
$\eta_{\infty} = 1/2$ and $c_{\eta} = -\beta/3$}\label{sec:discoveries}

The $C$-splitting law $\eta(N)$ was extracted as follows. For each $N$,
we fit $\eta(N)$ as a free scalar against the residual $C$-asymmetry
of \eqref{eq:c2} with $\Phi_{C}\equiv 1$, giving the five values
\[
\eta(3)=0.405,\quad \eta(4)=0.409,\quad \eta(5)=0.487,
\]
\[
\eta(6)=0.485,\quad \eta(8)=0.448.
\]
A one-parameter functional fit $\eta(N) = a + b/N^{2}$ returns
$a=0.486(\pm 0.014)$, $b=-0.764(\pm 0.04)$, both within $2$--$3\%$ of
the rationals $a=1/2$ and $b=-16/21$. The fit
$\eta(N) = 1/2 - 16/(21 N^{2})$ predicts
\[
\eta(6)\simeq 0.479\quad\text{vs}\quad 0.485\;\text{measured},
\]
i.e.\ $1.3\%$ agreement on the most precisely measured AT2021 point.

\paragraph*{Interpretation.} The asymptotic value $\eta_{\infty} = 1/2$
is the structural large-$N$ statement \emph{the $C=-1$ tube is on
average $50\%$ heavier than the $C=+1$ tube}. We do not commit to a
mechanism here; candidate routes include (i)~spin-1/2 weight in the
adjoint representation, (ii)~axial $\mathbb{Z}_{2}$ anomaly, and
(iii)~the ratio of complex to real representation dimensions.

\paragraph*{The $-\beta/3$ link.} That the leading $1/N^{2}$
correction is exactly $-\beta/3$ ties the $C$-sector to the
parity sector through a single rational. The factor of $3$
matches the three transverse dimensions of the flux tube, hinting at
a geometric (rather than dynamical) origin.

\section{Honest limitations}\label{sec:limits}

\textbf{(L1) Absolute scale.} The string tension $\sqrt{\sigzero}$
enters as a dimensional input. We do not predict it from
arithmetic; the no-go theorems of \cite{RemondiereNoGo2026} (in
particular T1 of that work, on dimensional transmutation) make this
explicit.

\textbf{(L2) $14\%$ residual.} The mean absolute residual on $78$
filtered channels is $\sim 14\%$. This is comparable to the published
AT2021 systematic-plus-statistical errors on individual masses and on
the large-$N$ continuum extrapolation, but is by no means competitive
with a precision lattice calculation. Lüscher finite-volume
corrections \cite{Luscher1986a,Luscher1986b}, applied on top of
\eqref{eq:c2}, recover at best $5\%$ further on our data, suggesting
that AT2021 published values are already extrapolated to infinite
volume to within a few percent.

\textbf{(L3) The $2^{+-}$ channels.} Five $2^{+-}$ channels (one per
$N$ in our sample) are systematically over-predicted by $50$--$60\%$ in
the closed form. These are known to be ditorelon-dominated
multi-glueball scattering states \cite{AthenodorouTeper2021}, mixing
into the lowest single-particle eigenvalue from above. The formula
\eqref{eq:c2}--\eqref{eq:phi-C} is therefore restricted to channels
that are not contaminated by such low-lying scattering states.

\textbf{(L4) Cross-group test pending.} Bennett \emph{et al.}
\cite{Bennett2021} have published an analogous lattice spectrum for
$\mathrm{Sp}(2N)$ with $N\in\{1,2,3\}$. A preliminary fit with the
same $K^{2}$ and a group-specific $F_{\mathrm{Sp}}(N)$ agrees to
$2.5\%$ at the ground state, confirming $K$ as group-universal;
the corresponding cross-group test of $c^{2}$ and of $\eta(N)$ is
ongoing.

\textbf{(L5) Lean status.} Only $K$, $F$, and $\xistar$ are PROVED in
Lean 4. The spectral factor $c^{2}$ and the discoveries
$\eta_{\infty}=1/2$, $c_{\eta}=-\beta/3$ are empirical
(\textsc{Tier 3, supported}).

\section{Falsifiable predictions}\label{sec:predictions}

The closed form \eqref{eq:c2}--\eqref{eq:phi-C} furnishes
parameter-free predictions for as-yet-unmeasured channels. Three of
them have been singled out below for near-term lattice falsification.

\begin{table}[t]
\centering
\caption{Falsifiable predictions for selected channels on the
$\mathrm{SU}(N)$ lattice. Predictions assume
$\sqrt{\sigzero}=440$~MeV for the $\mathrm{SU}(3)$ entries.}
\label{tab:predictions}
\begin{tabular}{llr}
\toprule
Quantity & Predicted & Test \\
\midrule
$m(1^{++})_{\mathrm{SU}(3)}$ &
$\simeq 2501$~MeV &
new measurement \\
$m(6^{++})_{\mathrm{SU}(3)}$ &
$\simeq 6.0\sqrt{\sigzero}\cdot J$ (Regge) &
new measurement \\
$\eta_{\mathrm{SU}(6)}$ &
$0.479$ &
$0.485$ \cite{AthenodorouTeper2021} \\
$\eta_{\mathrm{SU}(\infty)}$ &
$1/2$ exactly &
$\mathrm{SU}(10)$ extrapolation \\
\bottomrule
\end{tabular}
\end{table}

A measurement deviating from any of these by more than $5\sigma$
would falsify the corresponding sector of \eqref{eq:c2}. The
$\eta_{\mathrm{SU}(6)}$ prediction is the sharpest available
falsifier today: the closed form gives $0.479$, the published AT2021
value is $0.485$, in agreement at the $1.3\%$ level (consistent with
the AT2021 errors). A future $\mathrm{SU}(10)$ data set, by
contrast, would distinguish $1/2$ from competing
fits.

\section{Conclusion and roadmap}\label{sec:conclusion}

We have written down a closed form for the mass-gap ratio
$m^{2}/\sigzero$ of 4D pure $\mathrm{SU}(N)$ Yang--Mills in every
$J^{PC}(\mathrm{ex})$ channel at every $N\geq 3$, with all coefficients
rational and the universal transcendental $\sqrt{4\pi e/3}$ as the only
non-rational ingredient. The dimensional scale enters via $\sqrt{\sigzero}$.
On the $78$-channel AT2021 sample the formula matches the lattice data
at the $14\%$ level, comparable to the published lattice systematics.

The roadmap toward the Clay problem is as follows. (i) Independent
$\mathrm{Sp}(2N)$ tests \cite{Bennett2021} of $\eta(N)$, with the same
two rationals $\eta_{\infty}=1/2$ and $c_{\eta}=-\beta/3$. (ii) A
Lüscher-style $V$-correction analysis on a precision SU(3) lattice to
pin down the residual $5\%$. (iii) Lean-4 mechanization of the
spectral factor $c^{2}$, extending the existing $K, F, \xistar$ kernel.
(iv) An arithmetic derivation of $\beta = 16/7$ from a parity-genus
identity, currently the most opaque rational in
\eqref{eq:c2}--\eqref{eq:phi-C}.

\medskip
\noindent\emph{Acknowledgements.}
The author thanks the \texttt{crossed-cosmos} open notebook project for
the lattice data tabulations used in the cross-$N$ verification. All
numerical work was carried out in \texttt{PARI/GP 2.15}, with
cross-checks in \texttt{Python}. The Lean~4 proofs of
$K$, $F$, $\xistar$ are kernel-verified at commit
\verb|Crossed/Transport.lean|.

\medskip
\noindent\emph{Data availability.}
Lattice data are from \cite{AthenodorouTeper2021,Bennett2021};
worksheets are at the project repository.

\begin{thebibliography}{99}

\bibitem{JaffeWitten2000}
A.~Jaffe and E.~Witten,
\emph{Quantum Yang--Mills theory},
Clay Mathematics Institute Millennium Problem description (2000).

\bibitem{AthenodorouTeper2021}
A.~Athenodorou and M.~Teper,
\emph{$\mathrm{SU}(N)$ gauge theories in $3+1$ dimensions: glueball
spectra, string tensions and topology},
JHEP \textbf{12} (2021) 082, \texttt{arXiv:2106.00364}.

\bibitem{Bennett2021}
E.~Bennett \emph{et~al.},
\emph{Sp(2N) lattice gauge theories and extensions of the
Standard Model of particle physics},
EPJ Plus \textbf{139} (2024); \texttt{arXiv:2103.00665}.

\bibitem{Morningstar1999}
C.~J.~Morningstar and M.~J.~Peardon,
\emph{The glueball spectrum from an anisotropic lattice study},
Phys.~Rev.~D \textbf{60} (1999) 034509;
\texttt{arXiv:hep-lat/9901004}.

\bibitem{LucaciuTeperWenger2004}
B.~Lucini, M.~Teper, and U.~Wenger,
\emph{Glueballs and $k$-strings in $\mathrm{SU}(N)$ gauge theories:
calculations with improved operators},
JHEP \textbf{06} (2004) 012; \texttt{arXiv:hep-lat/0404008}.

\bibitem{Jaynes1957a}
E.~T.~Jaynes,
\emph{Information theory and statistical mechanics},
Phys.~Rev.~\textbf{106} (1957) 620.

\bibitem{Jaynes1957b}
E.~T.~Jaynes,
\emph{Information theory and statistical mechanics. II.},
Phys.~Rev.~\textbf{108} (1957) 171.

\bibitem{ElstrodtGrunewaldMennicke1998}
J.~Elstrodt, F.~Grunewald, and J.~Mennicke,
\emph{Groups acting on hyperbolic space: harmonic analysis and number
theory}, Springer Monographs in Mathematics (1998).

\bibitem{Sarnak1983}
P.~Sarnak,
\emph{The arithmetic and geometry of some hyperbolic three-manifolds},
Acta Math.~\textbf{151} (1983) 253--295.

\bibitem{RemondiereXiStarCR2026}
K.~R\'emondi\`ere,
\emph{A topological invariant of Bianchi 3-orbifolds from the identity
term of the Selberg pretrace formula},
preprint, Comptes Rendus de l'Acad\'emie des Sciences (May 2026).

\bibitem{tHooft1974}
G.~'t~Hooft,
\emph{A planar diagram theory for strong interactions},
Nucl.~Phys.~B \textbf{72} (1974) 461.

\bibitem{DijkgraafWitten1990}
R.~Dijkgraaf and E.~Witten,
\emph{Topological gauge theories and group cohomology},
Commun.~Math.~Phys.~\textbf{129} (1990) 393--429.

\bibitem{LuscherWeisz2002}
M.~L\"uscher and P.~Weisz,
\emph{Quark confinement and the bosonic string},
JHEP \textbf{07} (2002) 049; \texttt{arXiv:hep-lat/0207003}.

\bibitem{LuscherWeisz2004}
M.~L\"uscher and P.~Weisz,
\emph{String excitation energies in $\mathrm{SU}(N)$ gauge theories
beyond the free-string approximation},
JHEP \textbf{07} (2004) 014; \texttt{arXiv:hep-th/0406205}.

\bibitem{Luscher1986a}
M.~L\"uscher,
\emph{Volume dependence of the energy spectrum in massive
quantum field theories. I.~Stable particle states},
Commun.~Math.~Phys.~\textbf{104} (1986) 177.

\bibitem{Luscher1986b}
M.~L\"uscher,
\emph{Volume dependence of the energy spectrum in massive
quantum field theories. II.~Scattering states},
Commun.~Math.~Phys.~\textbf{105} (1986) 153.

\bibitem{Maldacena1998}
J.~M.~Maldacena,
\emph{The large-$N$ limit of superconformal field theories and
supergravity},
Adv.~Theor.~Math.~Phys.~\textbf{2} (1998) 231--252;
\texttt{arXiv:hep-th/9711200}.

\bibitem{ChandraChevyrevHairerShen2024}
A.~Chandra, I.~Chevyrev, M.~Hairer, and H.~Shen,
\emph{Stochastic quantisation of Yang--Mills--Higgs in 3D},
Invent.~Math.~(2024); \texttt{arXiv:2201.03487}.

\bibitem{RemondiereNoGo2026}
K.~R\'emondi\`ere,
\emph{No-go theorems for the arithmetic determination of the
Yang--Mills mass gap absolute scale},
preprint, Physical Review Letters submission (May 2026).

\end{thebibliography}

\end{document}
